\subsection{THETA3-C30: simultaneous Fermat-base changes add only CM Albanese factors}\label{C30}

C28 rules out changing the base of one Schoen fibration.
There is also an exact restriction on the natural attempt
to change two bases simultaneously and use a new Prym
factor.  In this construction every additional Albanese
factor is a quotient of a power of the Fermat-quintic
Jacobian.  Thus no additional copy of the C25 target
appears, and its full C27 quartic is still zero.

\paragraph{Statement for the actual two-fibration cover.}
Let $C=C_5$ be the Fermat quintic
\[
                 U^5+V^5=W^5,
\]
and use the two morphisms
\[
 Y\longrightarrow C\times C,\qquad
 (y_{ij})\longmapsto
 \bigl((y_{13}:y_{01}:y_{03}),
       (y_{23}:y_{02}:y_{03})\bigr)
                                                        \tag{C30a}
\]
on the smooth Schoen ball quotient.  These are the
Fermat projections in Schoen, printed p.~133; C12
computed the generic degree of their product to be five.
Let $f_i:D_i\to C$ be connected finite \'etale covers
of degrees $d_i$, and set
\[
               X=Y\times_{C\times C}(D_1\times D_2).
                                                        \tag{C30b}
\]
Then $X$ is connected, smooth and projective, and is a
finite \'etale arithmetic ball cover of $Y$.  There is
an isogeny
\[
 \operatorname{Alb}(X)\ \sim\ J(D_1)\times J(D_2)\times B,
                                                        \tag{C30c}
\]
where $B$ is a quotient of
$J(C)^{75d_1d_2}$.  In particular $B$ is CM.
The isogeny is compatible, up to isogeny, with the
Albanese map induced by $X\to D_1\times D_2$.

For the simple fourfold $A$ of C25, every morphism
$h:X\to A$ therefore satisfies
\[
 [N]h=b\circ\pi+a_0
                                                        \tag{C30d}
\]
for some positive integer $N$, some morphism
$b:D_1\times D_2\to A$, and a constant $a_0\in A$;
here $\pi$ is the displayed map from $X$ to the product.
Consequently
\[
                    F_X=0
                                                        \tag{C30e}
\]
for the entire quartic in C27.  This includes all
Albanese maps on these covers, not only sums of the two
initially supplied curve maps.

\paragraph{The exact cyclic model and its root line.}
Let $H=\mathcal O_C(1)$, and denote the pulled-back
sections of $f_i^*H$ by $U_i,V_i,W_i$.  On
$P=D_1\times D_2$ put
\[
 H_i=f_i^*H,\qquad
 L=\operatorname{pr}_1^*H_1\otimes
                  \operatorname{pr}_2^*H_2.
\]
The degrees of the $H_i$ are $5d_i>0$.  Define
\[
 s=U_1^5W_2^5-W_1^5U_2^5\in H^0(P,L^5),
 \qquad
 Z=\operatorname{Spec}_P
       \left(\bigoplus_{j=0}^4L^{-j}\right),
                                                        \tag{C30f}
\]
where multiplication uses $s$ to identify the fifth
power of the root with $s$.  Thus $Z\to P$ is finite
flat of degree five, and its direct image of the
structure sheaf is precisely the displayed sum.

On $W_1W_2\ne0$, write $t_i=(U_i/W_i)^5$.
The function field is obtained by adjoining $w$ with
\[
                         w^5=t_1-t_2.
                                                        \tag{C30g}
\]
This is exactly the function field of (C30b): in the
original equations it is the remaining coordinate
$w=y_{12}/y_{03}$.  The branch divisor has smooth
points at which $t_1-t_2$ has a simple zero.  Its
pullback to $P$ still has a simple zero, because both
base changes are \'etale.  Hence this function is not
a fifth power.  Since five is prime and the ground
field contains its roots of unity, the cyclic field
extension has degree five.  In particular $Z$ is
integral.  The same calculation proves connectedness
of (C30b): every component of its finite \'etale map
to connected $Y$ would dominate $Y$, while its generic
algebra over $P$ is the single degree-five field just
computed.  Smoothness, projectivity and the arithmetic
ball property follow from its finite \'etale map to $Y$.

\paragraph{Every singularity is an ordinary Fermat cone.}
The map
\[
          t:C\longrightarrow\mathbf P^1,
                    t=U^5/W^5
\]
has degree $25$, with five points of ramification
index five over each of $0,1,\infty$, and no other
ramification.  Away from these three values, $t$ is
an \'etale coordinate.  Thus the divisor $s=0$ is
smooth there.  At $t=0$ use the local parameter
$x=U/W$, so that $t=x^5$.  At $t=1$ use
$x=V/W$, so that $t=1-x^5$.  At infinity use
$x=W/U$, so that $t=x^{-5}$.  These are actual local
parameters and remain so after the \'etale maps $f_i$.

At a pair of points above the same one of
$0,1,\infty$, the section in (C30f), in the indicated
local trivializations, is $x_1^5-x_2^5$ or its
negative.  At infinity this follows by trivializing
the two factors with $U_1,U_2$, which gives
$s/(U_1^5U_2^5)=x_2^5-x_1^5$.
Consequently the local cyclic surface is
\[
                       z^5=x_1^5-x_2^5,
                                                        \tag{C30h}
\]
after changing one coordinate by a fifth root of $-1$
when needed.  These are all its singular points.
There are exactly
\[
             3(5d_1)(5d_2)=75d_1d_2
                                                        \tag{C30i}
\]
of them.  At a point with just one $W_i=0$, the other
nonzero, the section in the relevant trivialization
is a unit, so the cyclic cover is \'etale and smooth.
This checks the points omitted by the affine equation
(C30g).

The surface is a hypersurface over the smooth surface
$P$, hence Cohen--Macaulay, and it is nonsingular in
codimension one.  It is therefore normal by Serre's
criterion, EGA~IV$_2$, Theorem~5.8.6, printed p.~108
(also Stacks, Lemma~10.157.4, tag \texttt{031S}).
Blowing up the vertex of (C30h) resolves it in one
step.  Indeed its exceptional divisor is the smooth
plane curve $z^5=x_1^5-x_2^5$, and in each standard
blow-up chart the strict transform is the corresponding
affine chart of that smooth curve times the radial
coordinate.  Thus the exceptional divisor is a copy
of $C$ with normal bundle $\mathcal O_C(-1)$.
Let
\[
                       r:\widetilde Z\longrightarrow Z
\]
be the resulting smooth projective resolution.  Its
exceptional divisor is a disjoint union of the
$75d_1d_2$ Fermat curves in (C30i).

The surfaces $X$ and $\widetilde Z$ are smooth
projective models of the same function field.  Their
Albanese varieties, maps to abelian varieties and
degree-four pullback integrals are consequently
identified birationally.  One may use Milne,
\emph{Abelian Varieties}, Theorem~3.2, printed p.~17:
a rational map from a nonsingular variety to an
abelian variety extends everywhere.  Applying it in
both directions gives the Albanese identification.
For the integrals, use a common smooth resolution;
each birational map has degree one.

\paragraph{No new coherent first cohomology on the finite model.}
For $1\leq j\leq4$, each $H_i^{-j}$ has negative
degree, so $H^0(D_i,H_i^{-j})=0$.  The coherent
K\"unneth formula therefore gives
\[
 \begin{aligned}
 H^1(P,L^{-j})={}&
 H^1(D_1,H_1^{-j})\otimes H^0(D_2,H_2^{-j})\\
 &\quad\oplus
 H^0(D_1,H_1^{-j})\otimes H^1(D_2,H_2^{-j})=0.
 \end{aligned}
                                                        \tag{C30j}
\]
In this case the formula follows directly by taking
affine \v Cech covers of the two curves: the product
cover computes the tensor product of the two
complexes, and the K\"unneth decomposition is ordinary
linear algebra over the ground field.
Finite pushforward and (C30f) now imply
\[
 H^1(Z,\mathcal O_Z)=H^1(P,\mathcal O_P),\qquad
 h^1(Z,\mathcal O_Z)=g(D_1)+g(D_2).
                                                        \tag{C30k}
\]
The isomorphism is the pullback from $P$, not merely
an equality of dimensions.

\paragraph{Resolving the cones adds only Fermat-Jacobian factors.}
Let $V=\operatorname{Alb}(\widetilde Z)$, and let
$B\subset V$ be the sum of the images of the
Jacobians of all its exceptional curves.  Thus $B$
is a quotient of $J(C)^{75d_1d_2}$.  Put $Q=V/B$.
The composite $\widetilde Z\to Q$ is constant on
each exceptional curve, because its induced Jacobian
homomorphism has been killed.  It is therefore constant
on every fibre of $r$.

This map descends to a morphism $Z\to Q$.  Here is
a direct descent argument that does not presume
rational singularities.  Take the reduced image of
its graph in $Z\times Q$.  Projection to $Z$ is
proper and birational, and its geometric fibres
consist of one point.  It is finite by the proper
finite-fibre criterion (Stacks, Lemma~37.44.1,
tag \texttt{02LS}).  A finite birational morphism
to normal $Z$ is an isomorphism, since its local
coordinate algebras are integral subalgebras of
the same function field.  This proves the descent.

The image of $\widetilde Z\to Q$ generates $Q$,
so pullback gives an injection
\[
         H^1(Q,\mathcal O_Q)
                         \hookrightarrow H^1(\widetilde Z,
                                                      \mathcal O_{\widetilde Z}).
\]
For example this follows from the surjection
$\operatorname{Alb}(\widetilde Z)\to Q$ and the
usual first-cohomology identification for the smooth
Albanese map.  It can also be seen by pulling back
holomorphic one-forms: a finite sum of differences
of the image generates a dominant map to $Q$, so
a form vanishing on the image must be zero; complex
conjugation gives the assertion for $H^1(\mathcal O)$.
As the map factors through $Z$, its pullback factors
through $H^1(Z,\mathcal O_Z)$.  Thus (C30k) yields
\[
                      \dim Q\leq g(D_1)+g(D_2).
                                                        \tag{C30l}
\]

The morphism $\widetilde Z\to P$ is dominant, and
the product Abel map generates
$J(D_1)\times J(D_2)$.  Hence it induces a surjection
\[
                    V\longrightarrow J(D_1)\times J(D_2).
\]
Every exceptional curve maps to a point of $P$, so
this map kills $B$ and factors through a surjection
from $Q$.  Its target has dimension equal to the
right side of (C30l).  It follows that
\[
                       Q\longrightarrow J(D_1)\times J(D_2)
\]
is an isogeny.  Poincar\'e complete reducibility now
gives (C30c), compatibly with this quotient; a full
proof of the required complement construction is
Milne, \emph{Abelian Varieties}, Proposition~10.1,
printed p.~42.  In particular no vanishing of
$R^1r_*\mathcal O_{\widetilde Z}$, and no assertion
that the cone singularities are rational, has entered
the proof.

For completeness, $J(C)$ is a product up to isogeny
of three CM abelian surfaces.  Its deck group over
the $t$-line is $(\Z/5\Z)^2$.  A character is given
by three entries in $\mathbf F_5$ of sum zero,
describing the inertia at $0,1,\infty$.  The only
characters with nonzero first cohomology have all
three entries nonzero; there are $4^2-4=12$ of them,
in three rational orbits of size four.  The associated
cyclic quotients have genus two by
Riemann--Hurwitz.  Pullback and trace give their
Jacobian product isogenous to $J(C)$, exactly as in
C28.  Each genus-two factor has a unital
$\Q(\zeta_5)$ action, of degree four, and is CM.
Every quotient of a power of this product is CM by
complete reducibility.  This also recovers the CM
Fermat input in Schoen's Remark~(ii), printed p.~134.

\paragraph{The C25 quotient and the entire quartic.}
A nonzero homomorphism from one of these genus-two
factors to the simple fourfold $A$ would have to
surject, which dimensions prohibit.  Consequently
\[
                         \Hom(B,A)=0.
\]
Every homomorphism $V\to A$ therefore factors through
$Q$, and the displayed isogeny $Q\to J(D_1)\times
J(D_2)$ makes it factor through that product after
multiplication by one positive integer.  Translating
the image of the chosen base point gives exactly
(C30d).  All isogeny denominators are retained in $N$.

Let $w$ lie in the two-dimensional Weil space of $A$.
Translations act trivially on cohomology, and
$[N]^*w=N^4w$.  Since $\pi$ has generic degree five,
the projection formula gives
\[
 N^4\int_X h^*w
       =5\int_{D_1\times D_2}b^*w=0.
                                                        \tag{C30m}
\]
The final equality is C14: the image of a product
of two curves is a Lefschetz class, and
$\operatorname{NS}(A)_\Q=\Q\theta$ for the C25
target.  Its Weil projection is zero.  Thus every
actual map has zero pairing, which proves that the
full quartic in C27 vanishes, including on its finite
search grid.

If the two original covers are instead finite
\'etale covers of $C_0$, first normalize their fibre
products with $C_5\to C_0$.  Each connected component
is a finite \'etale cover of $C_5$.  Applying the
argument to each pair of components proves the same
conclusion for every component of the simultaneous
$C_0$-base-change construction.  The statement does
not cover arbitrary finite \'etale covers of the
resulting ball surfaces.  Such a cover need not
arise from two base curves, and its new first
cohomology is not controlled by (C30f).

\framework{This is a restriction obtained from actual
normal models, a specified resolution, and the
Albanese universal property.  The original equations
and exceptional CM curves are Schoen, pp.~133--134.
The exact normality criterion is EGA~IV$_2$~5.8.6,
p.~108; graph descent uses the finite-morphism
criterion above.  The Albanese and realization
viewpoint is Grothendieck, cote~44, scans~7--8.
Together with C14 it excludes the new-Prym attempt
in this particular two-base tower.  It does not
construct an exceptional cycle or exclude covers
whose Albanese factors lie outside this tower.}
